\documentclass{article}
\usepackage[T1]{fontenc}
\usepackage{lmodern}
\usepackage{graphicx}
\usepackage{amsmath,amssymb}
\usepackage[a4paper,margin=2cm]{geometry}
\usepackage[absolute,overlay]{textpos}
\setlength{\TPHorizModule}{1mm}
\setlength{\TPVertModule}{1mm}
\usepackage[indonesian]{babel}
\usepackage{hyperref}
\setlength{\emergencystretch}{2em}
% unit: Halaman 63

\begin{document}

% === Halaman 63 (llm) ===

\begin{itemize}
  \item Biasanya, plainteks dimanipulasi dulu dengan kunci (umumnya di-XOR-kan dengan kunci eksternal, atau dengan kunci putaran pertama), lalu hasilnya ditransformasikan dengan fungsi $f$ berulang kali.
\end{itemize}

\vspace{10mm}

\begin{align*}
C_0 &= P \oplus K_0 \\
C_i &= f(C_{i-1}, K_i) \quad , i = 1, 2, ..., r-1 \\
C &= C_r = C_{r-1} \oplus K_r
\end{align*}

\vspace{10mm}

\begin{itemize}
  \item Semakin banyak jumlah putaran, efek diffusi semakin bertambah, namun jumlah putaran yang besar dapat memuat komputasi menjadi tidak sangkil (efisien).
  \item Jumlah putaran harus dipertimbangkan untuk mempertemukan kriteria keamanan dan efisiensi.
\end{itemize}

\end{document}
